\documentclass[conference]{IEEEtran}
\IEEEoverridecommandlockouts

\usepackage{cite}
\usepackage{amsmath,amssymb,amsfonts}
\usepackage{graphicx}
\usepackage{textcomp}
\usepackage{url}

\begin{document}

\title{ }

\author{\IEEEauthorblockN{ }
\IEEEauthorblockA{\textit{ } \\
\textit{ }\\
\mbox{ }}
}

\maketitle

\begin{abstract}
Equipment registers in higher education are commonly held in spreadsheets
that impose no constraints on what is recorded, so that the same item is
entered more than once and in more than one form. Replacing such a
spreadsheet with a governed register requires those records to be
reconciled, and it is this migration step, rather than the construction
of the system itself, that determines whether the resulting register can
be trusted. This paper studies how far that reconciliation can be
automated. Two corpora are combined: an established entity resolution
benchmark supplies human labelled duplicate pairs, and United Kingdom
Contracts Finder notices supply free text procurement descriptions
already carrying Common Procurement Vocabulary codes. A parameterised
degradation model reintroduces the error classes that published data
lacks, and lexical, embedding based and language model approaches are
compared for duplicate detection and for taxonomy assignment. Results are
reported not as headline accuracy but as the share of a register that can
be migrated without human intervention while precision is held at a fixed
target, together with the residual manual effort and the monetary cost of
each approach. [TBC] The resulting pipeline is deployed inside an asset
management system delivered to the Faculty of Computing, Engineering and
Science at the University of South Wales.
\end{abstract}

\begin{IEEEkeywords}
asset management, entity resolution, record linkage, data quality, text
classification, open data
\end{IEEEkeywords}

\section{Introduction}

Higher education institutions hold research and teaching equipment worth
many millions of pounds, and the records describing that equipment are
very often kept in spreadsheets that were never designed to act as an
authoritative register. A spreadsheet offers no access control beyond
file permissions, no audit trail, no way of attaching a photograph or a
calibration certificate to a particular item, and no mechanism for
warning a technician that a pressure vessel is due for inspection. It
also imposes no constraints on what is typed into it. Two members of
staff recording the same microscope in the same week will produce two
rows that differ in spelling, abbreviation, capitalisation and column
alignment, and the file will accept both without complaint.

The Faculty of Computing, Engineering and Science (FCES) at the
University of South Wales is currently relocating a large volume of
equipment into a new building, Calon. The existing record of that
equipment is a static spreadsheet. The faculty requires a live system in
which any item can be reached by scanning a QR or barcode label, updated
in place, and associated with photographs, technical documentation,
service history, a position on a floor plan and the relevant health and
safety risk assessment. Access has to be tiered, so that most staff can
read the register while a smaller group is able to modify it.

Building such a system is a software engineering exercise with a well
understood solution space. Populating it is not. The migration step, in
which several thousand free text rows written by different people over
many years have to be reconciled into a consistent structured register,
is the point at which asset management projects most often lose
accuracy, and it is the step that determines whether the finished system
is trusted by the people who depend on it. Duplicates that survive
migration cause the register to overstate holdings and distort
replacement planning. Categories assigned inconsistently make the
register impossible to search or report on. An item placed in the wrong
category will inherit whatever hazard classification and servicing rule
that category carries, so a classification error propagates directly
into a safety control.

These consequences are not hypothetical. In 2024/25 employers in Great
Britain reported 59,219 non fatal injuries to employees under the
Reporting of Injuries, Diseases and Dangerous Occurrences Regulations,
and the Health and Safety Executive estimates the annual cost of
injuries and new cases of work related ill health arising from current
working conditions at \pounds22.9 billion \cite{hse2025}. Handling,
lifting or carrying accounted for 17 per cent of reported non fatal
injuries and being struck by a moving object for a further 10 per cent
\cite{hse2025}. In both categories the correct identification of a
machine and its associated risk assessment is a practical control rather
than an administrative formality. A register that cannot reliably say
which machine is which is a register that cannot support those controls.

This paper therefore treats the migration problem itself as the object
of study rather than as preparatory work. The question is not whether a
live asset register can be built, because it plainly can, but how much
of the transition from an uncontrolled spreadsheet to a governed
register can be automated, how accurate that automation is, and how much
human effort remains once it has been applied.

This study is carried out alongside the delivery of that system, and the
two are related rather than parallel. The system provides item records
with attached photographs and documents, a QR label per item resolving to
a persistent item URL, pin placement on an uploaded floor plan, links to
the governing risk assessment, service interval tracking with scheduled
reminders, and tiered access across administrator, technician and read
only roles. The migration pipeline studied here runs inside it, as a bulk
import in which records the pipeline resolves confidently enter the
register directly and the remainder enter a review queue. The system is therefore both the destination of the migration and the
setting in which its residual human cost accrues and is recorded, which
is what would let that cost be measured in deployment instead of
assumed.

Answering that question empirically requires records for which the
correct answers are known. Records held by FCES were not available for
this study, and published institutional catalogues, being curated for
public display, carry neither duplicate annotations nor a consistent
category label. The study therefore draws on two corpora with
complementary ground truth. Established entity resolution benchmarks
supply human labelled duplicate pairs together with published baselines
against which the matching results can be positioned. Notices published
through the United Kingdom Contracts Finder service supply several
hundred thousand free text procurement descriptions already carrying
Common Procurement Vocabulary codes, which serve as category labels at a
scale no manual annotation exercise could reach and which sit in the same
descriptive register as equipment records. A parameterised degradation
model then reintroduces the error classes absent from both, so that method performance can be observed as
record quality falls rather than at a single arbitrary operating point.

The contributions of this paper are as follows.

\begin{itemize}
\item A two corpus evaluation design that separates duplicate ground
truth from category ground truth, allowing an asset migration pipeline to
be measured end to end without access to a proprietary register.
\item A parameterised degradation model, extending the dirty data
protocol used in the entity matching literature, that reproduces the
error classes documented in manually maintained asset and procurement
data at controllable severity.
\item A comparative evaluation of lexical, embedding based and language
model approaches to duplicate detection and to taxonomy assignment,
reported with accuracy, latency and monetary cost.
\item A cost aware cascade in which inexpensive similarity measures
resolve confident matches and non matches while a language model
adjudicates only the ambiguous band, together with a measurement of how
much of the workload this removes from the expensive tier.
\item An operating point analysis reporting the proportion of a register
that can be migrated without human intervention at a fixed precision
target, the residual manual effort this implies, and an implementation
of the pipeline inside an asset management system delivered to FCES.
\end{itemize}

The remainder of the paper is organised as follows. Section II states
the research questions and the scope of the study. Section III reviews
related work. Section IV describes the corpora, the degradation model
and the experimental design. Section V presents the results. Section VI
discusses their interpretation, the integration of the pipeline into the
delivered system and the threats to validity. Section VII concludes.

\section{Research Questions and Scope}

The study is organised around three research questions.

\textbf{RQ1.} How accurately can duplicate and near duplicate equipment
records be identified, and how does that accuracy degrade as the quality
of the underlying records falls?

Duplicate detection is treated first because every later stage depends
on it. A register containing the same item twice will apply the same
servicing rule twice and count the same asset twice. RQ1 compares
normalised exact matching, character level lexical similarity, sentence
embedding similarity and a hybrid cascade, measuring each across a range
of injected error rates rather than at a single operating point, so that
the sensitivity of each method to record quality is observed rather than
assumed.

\textbf{RQ2.} How accurately can free text equipment descriptions be
assigned to a controlled procurement taxonomy, and what accuracy, cost
and latency trade offs exist between classical, embedding based and
language model approaches?

Consistent categorisation is what makes a register searchable,
reportable and capable of carrying category level hazard and servicing
rules. RQ2 evaluates assignment to the Common Procurement Vocabulary,
which every public sector procurement notice in the United Kingdom
already carries and which is published as linked open data across the
higher education sector. The vocabulary is hierarchical, and performance
is reported at the two-digit division level and the four-digit class
level separately, because the practical value of a coarse but reliable
assignment differs from that of a finer but less certain one. The
eight-digit leaf level is not evaluated: the label distribution in the
available data is too sparse at that depth for the result to carry
meaning, and the measured sparsity is reported in Section V.

\textbf{RQ3.} What proportion of a legacy register can be migrated
without human intervention at a fixed precision target, and what
residual manual effort does the remainder represent?

RQ3 converts the measurements produced by RQ1 and RQ2 into the quantity
that determines whether the approach is adopted in practice. A method
that is accurate on average but cannot be operated at high precision is
of little use to a faculty that has to sign off a register, because
every error it introduces has to be found again by hand. Reporting the
automated share at a precision floor, rather than a headline accuracy
figure, states the result in the terms in which the decision to adopt or
reject the pipeline is actually made.

Three boundaries are placed on the scope of the work. First, the study
addresses the migration of an existing register and not the prediction
of equipment failure. Servicing is treated as a scheduling and
notification problem driven by category and interval, and predictive
maintenance is left to further work. Second, the location model is a two
dimensional pin placed on an uploaded floor plan image rather than a
spatial or building information model, which matches the granularity at
which the faculty records location. Third, all corpora consist of
published open data and no personal or commercially sensitive
information is processed, so no ethical approval beyond the standard
faculty declaration was required.

\section{Related Work}

\subsection{Record Linkage and Entity Resolution}

The problem of deciding whether two records describe the same real world
entity has a formal treatment dating to the probabilistic framework of
Fellegi and Sunter \cite{fellegi1969}, which models the comparison of
record pairs as a decision problem with configurable error rates.
Subsequent work systematised the surrounding pipeline, in particular the
blocking step that reduces the quadratic comparison space to a tractable
candidate set, and the similarity measures applied within it
\cite{christen2012,papadakis2020}. The blocking schemes evaluated here are drawn from this line of work. This literature is largely evaluated on bibliographic and
commercial product data, and the question of how the methods behave on
institutional asset records has received little attention.

\subsection{Learned Approaches to Entity Matching}

Mudgal et al. \cite{mudgal2018} explored the design space of deep
learning architectures for entity matching and established the benchmark
suite on which most subsequent work reports. Their study also introduced
a protocol for constructing dirty variants of structured datasets by
relocating attribute values into a free text field, which is the closest
published analogue to the column misalignment observed in manually
maintained spreadsheets. Li et al. \cite{li2020ditto} showed that
pre-trained language models applied as cross encoders improve on these
architectures, and Narayan et al. \cite{narayan2022} reported that large
foundation models perform competitively on data wrangling tasks without
task specific training. These approaches assume a labelled training set
drawn from the same distribution as the deployment data, an assumption a faculty
migrating a single spreadsheet cannot satisfy.

\subsection{Cost and Efficiency of Language Model Pipelines}

Where a large model is invoked per record pair, the cost of a pipeline
scales with the number of decisions it makes rather than with its
accuracy. The cascade evaluated here follows from that
observation: an inexpensive similarity measure resolves the instances on
which it is confident, and the expensive model is invoked only within the
band where the cheap decision is unreliable. Cost is rarely reported
alongside accuracy in the entity matching literature, so what a given
level of matching performance costs to obtain remains largely open, and
that is the question this design is intended to answer.

\subsection{Text Classification against Procurement Taxonomies}

Assigning free text descriptions to a hierarchical controlled vocabulary
is a long standing text classification problem, addressed classically
with sparse lexical features and linear classifiers and more recently
with dense sentence representations \cite{reimers2019}. Two properties of
the present task are less commonly addressed. The vocabulary is
hierarchical, so a coarse assignment may be reliable where a fine one is
not, and the two levels carry different practical value. And the
assignment carries a safety consequence, since a category determines
which servicing and hazard rules an item inherits, so the precision at
which an assignment can be made matters more than its average accuracy.
Neither consideration arises when the same classification is performed
for procurement analytics, which is where most existing work sits.

\subsection{Equipment Cataloguing and Asset Management}

Equipment cataloguing in higher education has been addressed
operationally rather than analytically. Open source catalogue software
such as Kit-Catalogue \cite{kitcatalogue} and aggregation services such
as the Equipment Data Service \cite{jisc} provide the means to publish
and share institutional holdings, and asset management practice is
codified in ISO 55000 \cite{iso55000}. The data quality problems that
arise when records are entered without constraint are themselves long
studied \cite{rahm2000}, but that literature addresses cleaning as a
general activity rather than the construction of a register. Existing work addresses entity matching, duplicate detection,
data cleaning and equipment cataloguing as separate concerns. What is
absent is an account of the transition itself: how accurately the records
leaving an uncontrolled spreadsheet can be reconciled into a governed
register, at what cost, and with how much human effort left over. That
transition is where an asset management deployment succeeds or fails, and
it is the omission this study addresses.

\section{Methodology}

\subsection{Study Design}

The study separates two concerns that are usually conflated: whether a
method works against ground truth a human has verified, and whether it
continues to work when record quality degrades. Two corpora are used,
each supplying verified ground truth for one of the two tasks, and a
degradation model is applied to both so that the second concern can be
measured on each.

Duplicate ground truth is supplied by established entity resolution
benchmarks in which candidate pairs have been manually labelled. Category
ground truth is supplied by public procurement notices, which carry a
taxonomy code assigned by the purchasing authority at the time of
publication. Neither corpus is drawn from an institutional asset
register. That substitution is the principal limitation of the design,
and its consequences are set out in Section VI rather than left
implicit.

\subsection{Corpus A: Duplicate Ground Truth}

Matching accuracy is measured on benchmarks from the entity resolution
literature in which candidate pairs have been manually labelled as
matching or non matching \cite{mudgal2018,koepcke2010}. The text heavy
Abt-Buy benchmark is used, because at least one of its attributes
contains long free text, and that is the property distinguishing
equipment descriptions from purely structured records. The corpus comprises [TBC] labelled pairs at a positive rate of [TBC],
divided into the training, validation and test splits supplied with it;
the figures quoted are for the full labelled set, and results are
reported on the test split of [TBC] pairs.

Two considerations justify this choice. Product catalogue entries share
the structural characteristics of equipment records: a short designation
carrying brand and model information inside a single name string,
attached to a variable free text description written by another party. More importantly, published
results exist for these benchmarks, so the methods evaluated here can be
positioned against the literature rather than reported in isolation.
The limitation, that the domain is commercial rather than
institutional, is carried explicitly through to the threats to validity,
and is the reason the in domain comparison described below is reported
alongside the benchmark figures rather than in place of them.

\subsection{Corpus B: Category Ground Truth}

Category assignment is evaluated on notices published through the United
Kingdom Contracts Finder service and released in Open Contracting Data
Standard format under the Open Government Licence \cite{contractsfinder}.
Every notice carries at least one Common Procurement Vocabulary code
assigned by the purchasing authority, alongside a free text title and
description, which yields a large labelled set for the task of mapping a
description to a taxonomy code.

The corpus is filtered to the CPV divisions relevant to laboratory,
engineering and computing equipment, so that the label distribution
resembles that of a faculty register rather than that of public
procurement as a whole. The retained divisions are [TBC],
yielding [TBC] records; the resulting class distribution is reported in
Section V. The CPV hierarchy itself is
reconstructed from the code and description pairs carried by the notices,
with parents derived by truncation, so that division and class level
evaluation are supported from a single source without an external
vocabulary file.

This corpus plays a second role. Its records are drawn from the same
descriptive register as equipment entries, being short procurement
descriptions written by many different authorities, so degraded variants
of these records serve as the in domain material for the duplicate
detection experiments, and a filtered subset populates the delivered
system with demonstration data.

The deduplication experiments draw on a stratified sample of [TBC]
records rather than on the corpus in full. The bound is imposed by the
cost of adjudicating the cascade's band, which scales with the number of
candidate pairs, and it is set so that every pair in the band can be
adjudicated exactly at the severity levels reported rather than estimated
across more of them. Sampling reduces the number of records available for
duplicate detection; it does not alter their character, since the sample
is drawn from the same distribution and the classification experiments
continue to use the corpus in full.

Labels assigned by purchasing authorities are not error free, and the
publisher documents known quality issues in the released data
\cite{contractsfinder}. A stratified sample is therefore reviewed to estimate label noise, and
the rate is reported alongside the classification results. The sample
comprises [TBC] records, and the observed rate of disagreement with the
published code is [TBC].

Reporting the rate is not sufficient on its own, because label noise does
not act only as a ceiling on measured accuracy. Where the mislabelling is
systematic rather than clerical, a classifier trained on those labels
learns the error and is scored as correct for reproducing it, so measured
accuracy overstates the accuracy a user would experience. The direction
matters: a ceiling makes a result pessimistic, whereas a learned and
rewarded error makes it optimistic, and the second is the more dangerous
reading to leave uncorrected. Where the manual review identifies a
recurring mislabelling, the per-class results are examined for evidence
that the classifiers reproduce it, and the finding is reported with the
records concerned rather than inferred from the aggregate.

\subsection{Degradation Model}

Published data is cleaner than a working spreadsheet, so a degradation
model reintroduces the missing error classes. The model extends the
dirty data protocol of Mudgal et al. \cite{mudgal2018}, which relocates
attribute values into a free text field, with six further classes
drawn from the quality issues the Contracts Finder publisher documents in
the released data and from a manual audit of a sample of it: abbreviation
from a domain
lexicon, character level insertion, deletion, substitution and
transposition, inconsistent casing, whitespace perturbation, field
omission, and variation in the notation of units and voltages.

Each class is applied independently with a probability governed by a
single severity parameter, producing a family of corpora spanning
lightly to heavily degraded records. Duplicate pairs are generated by
producing a second degraded copy of a source record under an independent
draw from the same model, reproducing the case in which two members of
staff enter the same item without reference to one another. Near duplicate distractors are also generated by degrading records that
are similar but genuinely distinct. On the procurement corpus these are drawn from pairs
sharing a taxonomy class and exhibiting high title similarity while
arising from distinct procurement processes. Establishing distinctness is harder than it appears. The publisher mints
an identifier per notice and not per process, so two notices describing
one procurement carry different identifiers, and a rule keyed on them
admits genuine duplicates into the negative set. The published
procurement reference is a better signal but not a sufficient one, since
it frequently encodes the issuing organisation instead of the process,
and the buyer identifier does not reliably resolve to a single body.
Candidates are therefore filtered by reference and by buyer, and the
contamination remaining after each filtering stage is measured and
reported.

Residual contamination is not eliminated but estimated. A random sample
of the surviving candidates is adjudicated against a written protocol,
released with the code, and the rate is reported with a Wilson score
interval. Adjudication is performed by a language model and not by a
human annotator, which is a limitation and is treated as one: the
adjudicating model is not the model used in the matching cascade, so a
systematic bias would not appear on both sides of one comparison, but the
estimate carries no external validation and should be read as indicative
of magnitude rather than as a measured constant.

The measured rate is high enough to determine what this corpus can and
cannot support, and that determination is itself a result. Contamination
in a negative set leaves recall untouched, since recall is computed over
duplicates whose membership is known by construction, but it corrupts
precision directly: a matcher that correctly identifies a contaminated
pair as a duplicate is scored as having produced a false positive.
Precision and F1 are therefore not reported on the procurement corpus.
Recall, pair completeness and the behaviour of both under increasing
degradation are, and precision, F1 and the operating point of RQ3 are
carried by the benchmark corpus, whose labels were produced by other
parties and published.

This is stated as a finding and not as an omission. Constructing an in
domain hard negative set from published procurement data failed here, and
it failed for a reason that generalises: the publication model emits
several notices per procurement while supplying no identifier that
distinguishes a process from an organisation, so pairs that are textually
near identical and genuinely non-matching cannot be separated reliably
from pairs that are near identical because they describe one purchase.
Two successive rounds of rule refinement did not resolve it. Anyone
attempting the same construction on the same class of data should expect
the same obstacle. On the benchmark corpus distractors are drawn from pairs sharing
a leading token. Without them a
detector that cannot separate similar but distinct records will report
high recall for the wrong reason. The model parameters, the
abbreviation lexicon and the generation code are released so that every
corpus used here can be regenerated exactly.

\subsection{Duplicate Detection}

Four methods are compared in increasing order of cost. Normalised exact
matching, applied after case folding, whitespace collapsing and
punctuation removal, establishes the floor and quantifies how much of
the problem is trivially solvable. Lexical similarity uses character
level n-gram term frequency inverse document frequency vectors compared
by cosine similarity; character level features are preferred to
word level features because the dominant error classes operate within
words. Semantic similarity encodes the concatenated name and description
with a sentence embedding model \cite{reimers2019}, which is expected to
tolerate abbreviation and paraphrase better and character noise worse.

The fourth method is a cascade. Pairs above an upper similarity
threshold are accepted, pairs below a lower threshold are rejected, and
only the band between them is passed to a large language model for
pairwise adjudication. Both thresholds are selected on the development partition, at each
degradation severity independently, to satisfy the precision target of
RQ3 while minimising the size of the adjudicated band. Selecting once on
undegraded records and carrying the result across would not satisfy the
target on the data it is applied to, and is not done.

A threshold satisfies the target when the lower bound of a Wilson score
interval on its precision reaches the target, not when its point estimate
does. The distinction is not pedantic. A point estimate admits thresholds
whose precision rests on a handful of accepted pairs, and one estimated
from a single pair is not an estimate at all; such thresholds are
selected on the development partition and then fail on the held out data,
which is the behaviour a precision floor exists to prevent. Requiring
confidence rather than a favourable point estimate also scales the
evidence demanded with the strictness of the target, which a fixed
minimum on accepted pairs would not.

The target is not always attainable. Where no threshold on the base
similarity reaches the required precision, the upper threshold is
undefined, nothing is accepted without adjudication, and the band extends
to every pair not confidently rejected. This is reported, not avoided: the severity at which the target ceases to
be reachable by a base matcher alone is itself a result, and it bounds
what any pipeline of this shape can offer on records of that quality.

A distinction follows that the cascade's structure makes easy to miss.
The threshold guarantees the target only on the pairs it accepts without
adjudication. Pairs resolved by the language model carry no such
guarantee, since no threshold governs them, so the precision of the
cascade's combined output is not bounded below by the target and may fall
beneath it. Both quantities are therefore reported: the precision of the
auto-accepted portion, which the selection procedure constrains, and the
precision of the combined output, which it does not. Where the two
diverge the difference is attributable to the adjudicator, and the
severity at which the combined output first falls below the target is
reported as the point beyond which the pipeline cannot deliver the
guarantee it was designed around. The proportion of pairs reaching the language model, and the resulting
cost per thousand records, are reported alongside accuracy, because a
cascade that improves accuracy at unbounded cost is not a usable result.

The three methods above are evaluated across the full factorial of
degradation severities and repetitions on both corpora, since none of
them carries a marginal cost per decision. The cascade is evaluated on
the benchmark corpus only. Its result is a trade of precision against
cost, and precision is not reportable on the procurement corpus, so
adjudicating a band there would consume budget to measure nothing the
paper states. On the benchmark corpus each adjudication is a paid call,
so the cascade is evaluated at a reduced set of severity levels and a
single repetition, with every pair in the band adjudicated at each. This
trades the resolution of the cascade's degradation curve for exactness at
the points measured, which is the better trade when the alternative is an
estimate at every point. The severity levels retained span the range
measured for the other methods, so the cascade can be read against their
curves rather than in isolation.

All methods operate over candidate pairs produced by blocking rather
than the full comparison space. Three keying schemes are evaluated: sorted
character n-grams of the normalised title, which applies to both corpora;
a leading token key taking the first substantive token of the title,
which approximates a brand designation where the title carries one; and
the identifier of the publishing authority, which is available on the
procurement corpus only. Blocking is evaluated separately by pair completeness and reduction
ratio, since recall lost during blocking cannot be recovered downstream,
and the availability of each scheme differs between the two corpora,
which is reported alongside its performance. The scheme parameters are selected on the
benchmark development partition, subject to a stated floor on pair
completeness, and carried unchanged to the procurement corpus in keeping
with the transfer design; any loss of completeness that transfer incurs
is reported rather than recovered by refitting. Because completeness lost
in blocking bounds the recall of every method downstream of it, a second
configuration selected on the procurement corpus is also reported, so
that the comparison between matchers remains measurable in domain even
where the transferred configuration does not support it. The two are
reported separately and never combined: the first answers what transfers,
the second answers how the matchers rank.

\subsection{Category Assignment}

Three approaches are compared, one from each family and in ascending
order of cost per decision. A lexical baseline represents each record as
character level n-gram features and classifies with a linear support
vector machine. An embedding baseline encodes the record with a sentence
embedding model and classifies with a logistic regression head. A
retrieval augmented language model condition is given a candidate label
set shortlisted by embedding similarity, rather than the full vocabulary,
together with the nearest labelled examples from the development
partition supplied as in context examples. Shortlisting is a practical
necessity as well as a design choice, since presenting a taxonomy of this
size in a prompt is neither affordable nor useful.

The comparison asks whether a language model given the labels already
available in Corpus B outperforms a classifier trained on the same
labels, which is the question facing anyone deciding how to categorise a
register. The narrower question of what the in context examples
contribute, which would require a second language model condition
without them, is left to further work. The two classical approaches are
evaluated on the full test partition; the language model condition is
evaluated on a stratified sample of it, whose size is stated with the
results, and the classical figures are reported on that same sample
alongside their full-partition figures so that the sample can be seen to
be representative.

The two levels present different difficulties. At division level the
label set is small and its distribution only moderately skewed, so the
task is tractable but coarse. At class level the label set is an order of
magnitude larger and its tail is long: most classes carry too few
examples to be learned at all. Evaluation at class level is therefore restricted, on two grounds
reported apart because they are different problems. Labels holding fewer
than a stated minimum of training examples are excluded as unlearnable.
Records whose published code names a division and assigns no class are
excluded as well, since they carry no class level ground truth at all;
scoring them would ask whether a classifier can predict that a buyer was
vague, which is a property of publishing practice and not of the
equipment. The proportion each restriction removes is reported with the
results, so class level accuracy is read against the share of the
register it applies to, not against the whole. Records falling outside
the scored set are not discarded from the migration pipeline; they are
routed to review, which is where both an unlearnable category and an
absent one belong.

Macro averaged scores are the primary measure at both levels, weighted
averages are reported alongside them, and a per class breakdown is given
for the divisions carrying hazard or servicing rules in the delivered
system, because those are the classes in which an error has an
operational consequence.

\subsection{Partitioning}

Every corpus is partitioned before any tuning. A development partition
is used for threshold selection, prompt design and model fitting, and a
held out partition is used only for the final reported results. In
Corpus A the published train, validation and test splits are used
unchanged, so that results remain comparable with the literature. In
Corpus B notices are partitioned by publication period rather than at
random, so that near identical repeat notices from the same buyer do not
appear on both sides, and so that the split reflects the way a register
accumulates over time. Without these constraints the
reported figures would be optimistic for reasons unrelated to method
quality.

\subsection{Evaluation Metrics}

Duplicate detection is evaluated by precision, recall and F1 on the
benchmark corpus, reported separately at every severity level rather than
summarised, since the components behave differently as records degrade.
On the procurement corpus only recall is reported, for the reason given
in Section IV: its negative set carries measured contamination that
depresses precision independently of any matcher. Blocking is evaluated
by pair completeness and reduction ratio. Category assignment is
evaluated by macro and weighted F1 at both taxonomy levels, with
confusion matrices for the operationally significant classes.

The operating point of RQ3 is computed from the benchmark corpus
severity sweep, since it is defined at a precision target and precision
is measured there.

Operational measures are reported for every method: mean processing
latency per record, monetary cost per thousand records for any method invoking a language
model, costed at the published rate card recorded with the run, and the automated share, defined as the
proportion of records the pipeline resolves without human intervention
while holding precision at or above the target specified in RQ3.
Residual manual effort is reported as the count and proportion of records
the pipeline routes to review, and total effort is expressed as that
count multiplied by a per record handling time and not fixed by a single
measured constant. No handling time is measured here. One observer's
reading speed would not generalise to the staff who would perform the
curation, and quoting it as though it did would put a precision on the
effort estimate that the evidence does not support. A reader holding a
figure for their own institution obtains their own total; a reader
without one still has the review volume, which is the quantity the
automated share determines. The review queue in the delivered system
records the duration of each decision, so the figure becomes measurable
in deployment without being asserted in advance. The precision target is
[TBC].

\subsection{External Validation}

Results are checked in two ways. First, the error distribution produced
by the degradation model is compared against the manual audit of the
Contracts Finder sample, and agreement or divergence in the relative
frequency of each error class is reported. A degradation model that
overrepresents an error class the methods handle easily would flatter
every subsequent result, so this comparison is a precondition for
interpreting the rest.

Second, the pipeline is tuned on Corpus A and then evaluated, without
further tuning, on degraded records from Corpus B. The comparison is made
on recall and pair completeness, which the contamination in Corpus B's
negative set does not affect. The thresholds
selected on the Corpus A development partition are carried across
unchanged; no parameter is refitted on Corpus B, and the transfer is
reported as a single paired comparison rather than as two independent
results. Corpus A carries human
labels but a commercial domain; Corpus B carries a domain closer to
institutional equipment but synthetic duplicate labels. Neither is a
substitute for a real asset register, and the difference between the two
figures is the best available estimate of how far the reported
performance would transfer. That difference is reported whatever its
size, and it is the first item in the threats to validity discussion.

\subsection{Integration into the Delivered System}

The pipeline is deployed inside the asset management system built for
FCES so that its behaviour can be observed under realistic use, with a
filtered subset of Corpus B standing in for the faculty register until
the real records become available. The
system provides item records with attached photographs and documents, a
QR label per item resolving to a persistent item URL, pin placement on
an uploaded floor plan image, service interval tracking with scheduled
reminders, tiered access comprising administrator, technician and read
only roles, and an audit log of changes.

Migration is exposed as a bulk import in which an uploaded spreadsheet
passes through the deduplication and classification stages and the
outcome is routed to one of two destinations. Records resolved above the
precision threshold are written directly to the register. Records below
it enter a review queue presenting the candidate decision, the competing
alternatives and the similarity evidence, so the residual effort reported
here corresponds to a task a curator can actually perform and not to an
abstract count of uncertain records. The queue records the duration of
each decision, which is what makes the per record handling time an
institution would need obtainable from its own use of the system.

\subsection{Reproducibility}

Collection scripts, schema mappings, the degradation model, the
annotation protocol and experimental code are released in a public
repository. Corpora derived from third party
sources are distributed as record identifiers and reconstruction scripts
rather than as redistributed data, in accordance with each licence.
Model versions, prompt texts, embedding checkpoints and random seeds are
recorded with every run. The language model tier uses an open-weights
model, so it is reproducible independently of any single provider's
endpoint remaining available. The repository is available at [TBC].

\section{Results}

\section{Discussion}

\section{Conclusion}

\begin{thebibliography}{99}

\bibitem{fellegi1969}
I.~P. Fellegi and A.~B. Sunter, ``A theory for record linkage,'' \emph{J. Amer.
  Statist. Assoc.}, vol.~64, no.~328, pp. 1183--1210, 1969.

\bibitem{christen2012}
P.~Christen, \emph{Data Matching: Concepts and Techniques for Record Linkage,
  Entity Resolution, and Duplicate Detection}. Berlin, Germany: Springer, 2012.

\bibitem{papadakis2020}
G.~Papadakis, D.~Skoutas, E.~Thanos, and T.~Palpanas, ``Blocking and filtering
  techniques for entity resolution: A survey,'' \emph{ACM Comput. Surv.},
  vol.~53, no.~2, 2020.

\bibitem{rahm2000}
E.~Rahm and H.~H. Do, ``Data cleaning: Problems and current approaches,''
  \emph{IEEE Data Eng. Bull.}, vol.~23, no.~4, pp. 3--13, 2000.

\bibitem{mudgal2018}
S.~Mudgal \emph{et~al.}, ``Deep learning for entity matching: A design space
  exploration,'' in \emph{Proc. ACM SIGMOD Int. Conf. Manage. Data}, 2018, pp.
  19--34.

\bibitem{koepcke2010}
H.~K\"{o}pcke, A.~Thor, and E.~Rahm, ``Evaluation of entity resolution
  approaches on real-world match problems,'' \emph{Proc. VLDB Endowment},
  vol.~3, no.~1, pp. 484--493, 2010.

\bibitem{li2020ditto}
Y.~Li, J.~Li, Y.~Suhara, A.~Doan, and W.-C. Tan, ``Deep entity matching with
  pre-trained language models,'' \emph{Proc. VLDB Endowment}, vol.~14, no.~1,
  pp. 50--60, 2020.

\bibitem{narayan2022}
A.~Narayan, I.~Chami, L.~Orr, and C.~R\'{e}, ``Can foundation models wrangle
  your data?'' \emph{Proc. VLDB Endowment}, vol.~16, no.~4, pp. 738--746, 2022.

\bibitem{reimers2019}
N.~Reimers and I.~Gurevych, ``Sentence-BERT: Sentence embeddings using Siamese
  BERT-networks,'' in \emph{Proc. Conf. Empirical Methods Natural Lang.
  Process.}, 2019, pp. 3982--3992.

\bibitem{iso55000}
\emph{Asset Management --- Overview, Principles and Terminology}, ISO Standard
  55000:2014, 2014.

\bibitem{hse2025}
Health and Safety Executive, ``Health and safety statistics for Great Britain,
  2024/25,'' 2025. [Online]. Available: https://www.hse.gov.uk/statistics/

\bibitem{jisc}
Jisc, ``Equipment data service.'' [Online]. Available:
  https://equipment.data.ac.uk/

\bibitem{contractsfinder}
Cabinet Office, ``Contracts Finder: Open Contracting Data Standard release
  packages,'' Open Government Licence v3.0. [Online]. Available:
  https://www.gov.uk/contracts-finder

\bibitem{kitcatalogue}
Loughborough University, ``Kit-Catalogue: Open source equipment cataloguing
  software,'' licensed under GPLv3.

\end{thebibliography}

\end{document}